% Part: lambda-calculus
% Chapter: lambda-definability
% Section: lambda-definable-recursive

\documentclass[../../../include/open-logic-section]{subfiles}

\begin{document}

\olfileid{lam}{ldf}{ldr}
\olsection{الدوال \usetoken{S}{lambda definable} عودية}

ليست جميع الدوال العودية الجزئية !!{lambda definable} فحسب، بل يصح
العكس أيضًا. أي إن جميع الدوال~!!{lambda definable} عودية جزئية.

\begin{thm}
  \ollabel{thm:lambda-computable} إذا كانت دالة جزئية~$f$
  !!{lambda definable}، فهي عودية جزئية.
\end{thm}

\begin{proof}
  نرسم البرهان رسمًا موجزًا فحسب. أولًا نعرّف ترقيمًا حسابيًا لحدود
  $\lambd$؛ أي نسند إلى حدود~$\lambd$ أعداد غودل منهجيًا، باستعمال
  الترميز المعتاد للمتتاليات بقوى الأعداد الأولية. ثم نعرّف دالة عودية
  جزئية~$\fn{normalize}(t)$ تأخذ عدد غودل~$t$ لحد لامبدا وسيطًا، وتعيد
  عدد غودل لصورته السوية إن كانت له صورة، وتكون غير معرّفة في غير
  ذلك. ثم نعرّف دالتين عوديتين جزئيتين~$\fn{toChurch}$ و
  $\fn{fromChurch}$ تحولان الأعداد الطبيعية إلى أعداد غودل لأعداد
  تشيرش الموافقة ومنها.

  وباستعمال هذه الدوال العودية يمكننا تعريف~$f$ دالة عودية جزئية.
  فهناك حد~$\lambd$~$F$ !!{lambda define}s~$f$. ولحساب
  $f(n_1, \dots, n_k)$، نحصل أولًا على أعداد غودل لأعداد تشيرش
  الموافقة باستعمال~$\fn{toChurch}(n_i)$، ثم نلحقها بـ$\Gn{F}$ لنحصل
  على عدد غودل للحد~$F \num{n_1}\dots\num{n_k}$. وبعد ذلك نطبق
  $\fn{normalize}$ في عدد غودل هذا. فإذا كانت~$f(n_1, \dots, n_k)$
  معرّفة، كان للحد~$F \num{n_1}\dots\num{n_k}$ صورة سوية (ولا بد أن
  تكون عدد تشيرش)، وإلا فلا صورة سوية له (ومن ثم تكون
  \[\fn{normalize}(\Gn{F\num{n_1}\dots\num{n_k}})\]
  غير معرّفة). وأخيرًا نطبق~$\fn{fromChurch}$ في عدد غودل للحد بعد
  تسويته.
\end{proof}

\end{document}
